% !TeX root = ../../Book3.tex
\section{Mittag–Leffler 定理}
\label{b3:sec:2004}

乘法构造指定零点，加法构造则指定极点附近的负幂项。把每个远处极点的低阶 Taylor 项扣掉，可以沿用同样的尾部控制思想。

\subsection{主要部分}
\label{b3:subsec:200401}

亚纯函数在极点 \(a\) 的 Laurent 负幂有限和 \(Q_a(z)=\sum_{k=1}^{m_a}c_{a,k}(z-a)^{-k}\) 称为\term[zhuyoubufen]{主要部分}[principal part]。剩余部分在 \(a\) 全纯。指定主要部分比只指定极点阶数强，因为它固定了每个负幂系数及留数。

\subsection{亚纯插值}
\label{b3:subsec:200402}

设 \(a_n\to\infty\)，给定主要部分 \(Q_n\)。选 \(R_n<|a_n|\)，且 \(R_n\to\infty\)，在圆盘 \(|z|\le R_n\) 上用 \(Q_n\) 关于零的 Taylor 多项式 \(T_n\) 逼近到误差 \(2^{-n}\)。这可行，因为 \(Q_n\) 在半径 \(|a_n|\) 内全纯，Taylor 级数在更小闭圆盘一致收敛。

于是 \(\sum(Q_n-T_n)\) 在避开极点的每个紧集上正规收敛。每个指定极点附近，除该项之外的尾部全纯，故其主要部分不被更改。

\subsection{Mittag–Leffler 定理}
\label{b3:subsec:200403}

\begin{theorem}[Mittag--Leffler：整平面版本]\label{b3:thm:20-mittag-leffler}
任意无有限聚点的极点集及其有限主要部分，都由一个复平面亚纯函数实现。两个实现函数之差为整函数。
\end{theorem}
\begin{proof}
除去可能在原点的有限主要部分后，用上一小节的 Taylor 扣除构造并相加。局部正规收敛给亚纯性，极点邻域中的有限项分离给准确主要部分。两解的差在每个极点负幂全部抵消，因此所有极点可去，成为整函数。
\end{proof}
一般平面区域上的相应定理需要把多项式扣除改成适应补集的有理逼近；本节给出的完整构造范围是整平面，不把这一证明无说明地扩展到任意区域。

\subsection{与 Weierstrass 定理的对偶性}
\label{b3:subsec:200404}

Weierstrass 定理留下无零点的乘法自由度 \(e^g\)，Mittag--Leffler 定理留下整函数的加法自由度。两者还能组合给整插值：取简单零点为 \(a_n\) 的整函数 \(P\)，给主要部分 \(b_n/(P'(a_n)(z-a_n))\)，构造亚纯函数 \(G\)，则 \(F=PG\) 整且 \(F(a_n)=b_n\)。局部乘法直接核验值，任意增长的插值数据也不破坏离散点上的存在性。


主要部分构造也有误差与有效范围的区别。对每个 \(Q_n\) 减去有限 Taylor 多项式，减去的项不引入新极点，并且改变的是整函数自由度。若只减常数项，在近原点的紧集上尾部可能仍不收敛；所需次数应根据该项大小选择，而不是统一宣称一次修正足够。

\ProseParagraphBreak

整插值还可以指定有限阶导数。设在 \(a_n\) 希望 \(F\) 与给定多项式芽 \(J_n\) 的差至少有 \(m_n\) 阶零点。取 \(P\) 在 \(a_n\) 有恰好 \(m_n\) 阶零点，令 \(Q_n\) 是 \(J_n/P\) 的主要部分。由 Mittag--Leffler 构造 \(G\)，则 \(G-J_n/P\) 在 \(a_n\) 全纯，故 \(F=PG\) 满足 \(F-J_n\) 被 \((z-a_n)^{m_n}\) 整除。负幂计算准确保留了全部所需导数；不需要把每阶导数当作独立极点重新构造。

\ProseParagraphBreak

这种插值没有控制全平面增长。若要求在整数上取值 \(e^{n^2}\) 且整函数阶小于二，就会与最大模的下界矛盾；若指定的值只是有界，也不能自动保证存在有界整函数，因为 Liouville 使有界整函数只能恒常。存在性定理容许任意数据，正是因为它没有许诺这些额外估计。


\ProseParagraphBreak

下面把 Taylor 扣除步骤写成一个可检查的例子。在正整数 \(n\) 指定二阶主要部分 \(c_n/(z-n)^2\)，其中 \(c_n\) 可以任意迅速增长。对 \(|z|<n\)，有
\[
\frac{c_n}{(z-n)^2}=\frac{c_n}{n^2}\sum_{k\ge0}(k+1)(z/n)^k.
\]
取 \(R_n=n/2\)，选择整数 \(d_n\) 使
\[
\frac{|c_n|}{n^2}\sum_{k>d_n}(k+1)2^{-k}<2^{-n}.
\]
有限截断 \(T_n\) 为整多项式，故 \(\sum_n(c_n/(z-n)^2-T_n(z))\) 在任意避开正整数的紧集上正规收敛。对任意固定 \(R\)，仅有限个 \(R_n<R\)；剩余项由 \(2^{-n}\) 一致控制。由于没有一次负幂项，构造出的函数在每个极点的留数仍为零，不能因为修正次数高而凭空出现留数。

\ProseParagraphBreak

若 \(c_n=1\)，可以直接用 \(\sum_{n\ge1}(z-n)^{-2}\)，因为远处尾部已经是 \(O_K(n^{-2})\)；若 \(c_n=n^2\)，直接尾项趋于一，减去常数项后尾项仍有 \(O_K(n^{-1})\)，还须减一次项。具体地，
\[
\frac{n^2}{(z-n)^2}-1-\frac{2z}{n}=O_K(n^{-2}).
\]
由此可见，扣除次数应从具体尾项的展开决定。一般定理提供可用算法，特殊数据则可用更低次数的修正。

\ProseParagraphBreak

Mittag--Leffler 构造的唯一性比乘积构造更容易表达：具有相同极点主要部分的两个亚纯函数之差在每个极点可去，因而为整函数。反过来，加上任意整函数都不改变主要部分。因此“主要部分决定函数”只有在另有增长、周期性或规范化时才可能成立。对有限个极点，所有主要部分之和已经给出有理函数，但一般解仍是这个有理函数加任意整函数；极点数有限并不使函数必为有理。

\ProseParagraphBreak

插值点必须离散。例如若在 \(1/n\) 指定值零，同时在原点指定值一，任何全纯函数的连续性已产生矛盾；即使原点也指定零，恒等定理仍迫使函数恒为零，不能再在另一点任意指定非零值。有限阶导数插值也没有绕过这个限制。所谓“任意数据”准确地指在没有有限聚点的点列上，每一点只提出有限个导数条件。
